% mazzi: 04
\chapter{Il tessuto epiteliale}

% ---------------------------------------------------------------------------
\section{I tessuti del corpo umano}

\textbf{Tessuti}: gruppi di cellule che cooperano, con simili caratteristiche morfologiche
$\rightarrow$ permettono una \textit{divisione del lavoro} fra le cellule dell'organismo. Nel
corpo umano si riconoscono \textbf{4 tipi}:
\begin{itemize}
  \item \textbf{epiteliale}: copre superfici esposte, delimita dotti e cavità interne, produce
    secreti ghiandolari;
  \item \textbf{connettivo}: riempie gli spazi interni, offre supporto strutturale, conserva energia;
  \item \textbf{muscolare}: si contrae per produrre un movimento attivo;
  \item \textbf{nervoso}: conduce impulsi elettrici e trasporta informazioni.
\end{itemize}

Organizzazione gerarchica: $molecole \rightarrow cellule \rightarrow tessuti \rightarrow organi
\rightarrow apparati$. I tessuti con funzioni speciali formano organi, che interagiscono negli
apparati.

% ---------------------------------------------------------------------------
\section{Caratteristiche generali}

\textbf{Tessuto epiteliale}: cellule morfologicamente differenziate, strettamente unite da
\textbf{giunzioni specializzate}, con scarsissima matrice extracellulare interposta.

\rubrica{Giunzioni cellulari}
\begin{itemize}
  \item \textbf{desmosomi}: ancoraggio meccanico tra cellule adiacenti;
  \item \textbf{giunzioni serrate} (\textit{tight junction}): sigillano lo spazio intercellulare,
    impediscono il passaggio di molecole tra le cellule;
  \item \textbf{giunzioni comunicanti} (\textit{gap junction}): comunicazione diretta tra cellule
    adiacenti.
\end{itemize}

\rubrica{Funzioni}
Protezione, scambio e lubrificazione di superfici interne ed esterne, tramite cellule
specializzate in assorbimento, secrezione, trasporto, scambio di gas, scivolamento e ricezione
sensoriale $\rightarrow$ la morfologia si differenzia in base alla funzione svolta.

\rubrica{Due categorie}
\begin{itemize}
  \item \textbf{epiteli di rivestimento}: ricoprono le superfici esterne ed interne del corpo;
  \item \textbf{epiteli ghiandolari}: specializzati nella produzione e nel rilascio di secreti.
\end{itemize}

% ---------------------------------------------------------------------------
\section{La membrana basale}

\textbf{Membrana basale}: presente alla base di tutti gli epiteli, formata da proteoglicani,
separa l'epitelio dal connettivo sottostante (tessuto connettivale reticolare lasso, ricco di
\textbf{collagene di tipo III}). È composta da due strati:
\begin{itemize}
  \item \textbf{lamina lucida}: strato più esterno, in prossimità della membrana plasmatica delle
    cellule epiteliali;
  \item \textbf{lamina densa}: strato più profondo, secreto dalle cellule del connettivo;
    contiene collagene di tipo IV (placche di ancoraggio) e di tipo VII (fibrille di ancoraggio).
\end{itemize}
Ricca di \textbf{molecole di adesione cellulare} (MAC) e \textbf{proteoglicani} (cemento
intercellulare) $\rightarrow$ assicurano l'ancoraggio delle cellule epiteliali.

% ---------------------------------------------------------------------------
\section{Polarità}

Ogni epitelio ha una \textbf{superficie esposta} (verso l'esterno o le cavità interne) e una
\textbf{superficie basale} (contatto con le strutture circostanti): differiscono per struttura e
funzione perché i componenti citoplasmatici non sono distribuiti uniformemente tra i due poli.

\rubrica{Strutture della superficie apicale}
\begin{itemize}
  \item \textbf{microvilli}: estroflessioni digitiformi $\rightarrow$ aumentano la superficie di
    assorbimento;
  \item \textbf{ciglia}: appendici mobili per il trasporto di muco e particelle;
  \item \textbf{stereociglia}: microvilli molto allungati (epididimo, orecchio interno).
\end{itemize}

% ---------------------------------------------------------------------------
\section{Classificazione degli epiteli di rivestimento}

\subsection{In base al numero di strati}
\begin{itemize}
  \item \textbf{monostratificato} (semplice): un unico strato, tutte le cellule a contatto con la
    lamina basale;
  \item \textbf{pseudostratificato}: tutte le cellule poggiano sulla lamina basale ma i nuclei
    sono a diverse altezze (apparenza di più strati);
  \item \textbf{pluristratificato} (composto): più strati sovrapposti; solo lo strato basale a
    contatto con la lamina basale.
\end{itemize}

\subsection{In base alla forma delle cellule}
\begin{itemize}
  \item \textbf{pavimentoso}: cellule piatte (larghezza e lunghezza $\gg$ spessore);
  \item \textbf{cubico}: forma approssimativamente cubica, nucleo centrale sferico;
  \item \textbf{cilindrico}: più alte che larghe, nucleo basale ovoidale.
\end{itemize}

\subsection{Epitelio cubico}
\begin{itemize}
  \item \textbf{semplice}: sede in ghiandole, dotti, tubuli renali, tiroide; funzioni limitate di
    protezione, secrezione e assorbimento;
  \item \textbf{stratificato}: ricopre alcuni dotti (raro); protezione, secrezione, assorbimento.
\end{itemize}

\subsection{Epitelio cilindrico}
\begin{itemize}
  \item \textbf{semplice}: sede in stomaco, intestino crasso, colecisti, tube uterine, dotti
    collettori del rene; protezione, secrezione, assorbimento;
  \item \textbf{stratificato}: piccole zone di faringe, epiglottide, ano, dotti mammari, ghiandola
    salivare, uretra; protezione.
\end{itemize}

\subsection{Epitelio cilindrico ciliato pseudostratificato}
Sede: cavità nasale, trachea e bronchi; porzioni dell'apparato riproduttivo maschile. Funzioni:
protezione e secrezione; le ciglia muovono il muco verso l'esterno $\rightarrow$ eliminazione di
particelle e agenti patogeni.

\subsection{Epitelio di transizione (urotelio)}
Sede: vescica urinaria, pelvi renale, ureteri. Permette espansione e ritorno dopo la distensione:
le cellule cambiano forma con lo stato di riempimento (più piatte a vescica piena, più arrotondate
a vescica vuota).

\subsection{Epitelio pavimentoso pluristratificato cheratinizzato --- l'epidermide}
Nell'\textbf{epidermide} le cellule superficiali perdono i nuclei e il citoplasma si riempie di
\textbf{cheratina} (proteina fibrosa resistente) $\rightarrow$ cellule non vitali =
\textbf{squame cornee} = barriera fisica contro traumi, disidratazione e agenti patogeni.

% ---------------------------------------------------------------------------
\section{Gli epiteli ghiandolari}

\subsection{Ghiandole esocrine ed endocrine}
\textbf{Epiteli ghiandolari}: specializzati nella produzione di secreti. In base alla destinazione:
\begin{itemize}
  \item \textbf{esocrini}: secreto riversato, tramite un \textbf{dotto escretore}, sulla
    superficie esterna o in cavità; la porzione secernente (\textbf{adenomero}) resta collegata
    all'epitelio dal dotto;
  \item \textbf{endocrini}: secreto (\textbf{ormoni}) riversato direttamente nei vasi sanguiferi;
    la porzione secernente si separa dalla superficie e viene diffusamente vascolarizzata.
\end{itemize}

\subsection{Classificazione delle ghiandole esocrine pluricellulari}
Quattro criteri:
\begin{enumerate}[label=\Alph*)]
  \item \textbf{in base alla posizione}:
    \begin{itemize}
      \item \textbf{intraparietali}: nella parete dell'organo:
        \begin{itemize}
          \item intraepiteliali: dentro lo strato epiteliale;
          \item esoepiteliali: sotto l'epitelio (coriali o sottomucose);
        \end{itemize}
      \item \textbf{extraparietali}: a distanza dall'organo (es.\ pancreas, fegato);
    \end{itemize}
  \item \textbf{in base alla struttura del dotto e della porzione secernente}:
    \begin{itemize}
      \item \textit{semplici} (dotto non ramificato): tubulare semplice (gh.\ intestinali, cripte
        di Lieberkühn), tubulo spirale semplice (gh.\ sudoripare merocrine), tubulo ramificata
        semplice (gh.\ gastriche, mucose di esofago/lingua/duodeno), alveolare (acinosa) semplice
        (non nell'adulto), alveolare ramificata semplice (gh.\ sebacee);
      \item \textit{composte} (dotto ramificato): tubulare composta (gh.\ mucose della bocca,
        bulbouretrali, tubuli seminiferi), alveolare (acinosa) composta (gh.\ mammaria),
        tubulo-alveolare composta (gh.\ salivari, tratto respiratorio, pancreas);
    \end{itemize}
  \item \textbf{in base al meccanismo di secrezione}:
    \begin{itemize}
      \item \textbf{merocrine}: rilascio per \textbf{esocitosi}, senza perdita di citoplasma (il
        più comune, es.\ gh.\ salivari);
      \item \textbf{apocrine}: secreto accumulato all'apice, poi distaccato; la cellula viene
        ripristinata (4 stadi: secrezione $\rightarrow$ decomposizione apicale $\rightarrow$
        ripristino $\rightarrow$ nuovo ciclo); es.\ gh.\ mammaria;
      \item \textbf{olocrine}: le cellule accumulano fino a ``scoppiare'', rilasciando l'intero
        citoplasma; rimpiazzate per mitosi delle staminali basali (es.\ gh.\ sebacee);
    \end{itemize}
  \item \textbf{in base al tipo di secreto}:
    \begin{itemize}
      \item \textbf{sierose}: secrezione acquosa ricca di enzimi (es.\ amilasi salivare);
      \item \textbf{mucose}: secernono \textbf{mucine} (glicoproteine) che assorbono acqua
        $\rightarrow$ muco viscoso;
      \item \textbf{miste}: diversi tipi di cellule ghiandolari, sia sierosa che mucosa.
    \end{itemize}
\end{enumerate}

% ---------------------------------------------------------------------------
\section{Principali funzioni del tessuto epiteliale}
\begin{enumerate}
  \item \textbf{protezione fisica}: barriera meccanica contro traumi, microrganismi, disidratazione;
  \item \textbf{regolazione della permeabilità}: controllo selettivo del passaggio di sostanze tra
    i compartimenti;
  \item \textbf{sensibilità}: ricezione di stimoli dall'ambiente esterno e interno;
  \item \textbf{produzione di secreti specializzati}: enzimi, ormoni, muco e altre molecole
    biologicamente attive.
\end{enumerate}
